\lyseis{Β3.3}{Παραλληλόγραμμα και τραπέζια}
\sol{1} α) ορθογώνιο β) ρόμβος γ) τετράγωνο.
\sol{2} α) Σωστό. β) Λάθος. γ) Σωστό. δ) Λάθος.
\sol{3} Ρόμβος: $4\cdot6{,}5=26$ cm. Τετράγωνο: $4\cdot4=16$ cm.
\sol{4} Ελεύθερη κατασκευή.
\sol{5} Ελεύθερη κατασκευή. Το ύψος είναι κάθετο στις βάσεις.
\sol{6} Όλες οι πλευρές είναι ίσες, άρα ρόμβο (και ειδικά τετράγωνο, αν οι γωνίες είναι ορθές).
\sol{7} Δύο ισοσκελή τρίγωνα, αυτά που έχουν βάσεις τις βάσεις του τραπεζίου και κορυφή το σημείο τομής των διαγωνίων.
\sol{8} Τετράγωνο. Περίμετρος: $4\cdot8=32$ μ.
\sol{9} $5\cdot(4\cdot2)=40$ μ.
\sol{10} $120+80+50+50=300$ cm $=3$ μ.
\sol{11} $120:10=12$ ρόμβοι.
\sol{12} Ελεύθερη κατασκευή.
\sol{13} Τα τετράγωνα.
\sol{14} Το ορθογώνιο όταν όλες οι πλευρές είναι ίσες, ο ρόμβος όταν όλες οι γωνίες είναι ορθές.
\sol{15} Έχει μόνο δύο παράλληλες πλευρές, ενώ το παραλληλόγραμμο έχει και τα δύο ζεύγη.
\sol{16} Τετράπλευρο $\to$ παραλληλόγραμμο $\to$ ορθογώνιο και ρόμβος $\to$ τετράγωνο. Το τραπέζιο είναι ξεχωριστά κάτω από το τετράπλευρο.
\sol{17} Ελεύθερη απάντηση.
\sol{18} Ελεύθερη κατασκευή.
